\section{Experiments}
\label{sec:experiments}

\subsection{Experimental Setup}
\label{sec:exp_setup}

\noindent\textbf{Datasets.}
We evaluate on five standard CIL benchmarks:
\textbf{ImageNet-R}~\cite{hendrycks2021manyfaces} (200 classes, artistic renditions),
\textbf{DomainNet}~\cite{peng2019domainnet} (200 classes across six domains),
\textbf{OmniBenchmark}~\cite{zhang2022omnibenchmark} (300 classes aggregated from multiple benchmarks),
\textbf{CIFAR-100}~\cite{krizhevsky2009learning} (100 classes),
and \textbf{CUB-200}~\cite{wah2011cub200} (200 fine-grained bird species).
For each dataset, we evaluate under two class-incremental settings by varying the number of classes per task, resulting in approximately 10-step and 20-step protocols per dataset. All class orders are randomized with fixed seeds following~\cite{zhou2025revisiting}.

\noindent\textbf{Backbone.}
All methods use a ViT-B/16~\cite{dosovitskiy2021vit} pretrained on ImageNet-21K~\cite{ridnik2021imagenet21k} as the frozen feature extractor, with image resolution $224 \times 224$.

\noindent\textbf{Metrics.}
We report two standard CIL metrics~\cite{van2019three}: the \textbf{average incremental accuracy} $\bar{A} = \frac{1}{T}\sum_{t=1}^{T} A_t$, where $A_t$ is the test accuracy on all classes seen up to task $t$, and the \textbf{final accuracy} $A_T$ after the last task. Both reflect the model's ability to classify across the full accumulated label space without task identity.

\noindent\textbf{Compared methods.}
We compare against three families of exemplar-free CIL methods. \emph{Prompt-based:} L2P~\cite{wang2022learning}, DualPrompt~\cite{wang2022dualprompt}, CODA-Prompt~\cite{smith2023coda}. \emph{Adapter and fine-tuning based:} APER~\cite{zhou2025revisiting}, EASE~\cite{ease2024}, SLCA~\cite{slca2023}, SSIAT~\cite{tan2023ssiat}, CoFiMA~\cite{cofima2024}. \emph{Prototype and expert-based:} FeCAM~\cite{fecam2023}, RanPAC~\cite{ranpac2023}, ACIL~\cite{acil2022}, TUNA~\cite{wang2025integrating}, MOS~\cite{mos2025}, MiN~\cite{wang2024min}. For all methods, we use the officially released code and follow their recommended hyperparameters. Methods marked with $^*$ in tables indicate reproduction from source code under our unified evaluation protocol.

\noindent\textbf{Implementation details.}
SPIE is implemented in PyTorch. All experiments are conducted on NVIDIA A100 (80G) GPUs. The backbone is frozen throughout. For each task, the class prototypes are trained for 20 epochs with SGD (momentum $0.9$, cosine annealing, initial learning rate $0.02$). Each expert is trained for 35--40 epochs with the same optimizer and learning rate $0.02$--$0.03$. The CosFace loss uses scale $s=20$ and margin $m=0$. Classifier alignment runs for 30 epochs with learning rate $0.005$. The low-rank covariance rank is set to $r=16$ by default (ablations use $r \in \{4, 8, 16, 32\}$). Each expert uses $K=4$ expert tokens and adapter rank $r_v=256$. The posterior temperatures $\tau_t$ and $\tau_e$ are set to $1.0$ (we provide sensitivity analysis in Sec.~\ref{sec:ablation}). Batch size is $32$ for all experiments.

\subsection{Main Results}
\label{sec:main_results}

% Table 1: ImageNet-R, OmniBenchmark, DomainNet

\begin{table}[t]
\centering
\caption{
Performance comparison on ImageNet-R, OmniBenchmark, and DomainNet with ViT-B/16-IN21K as the backbone.
The best performance is shown in bold, and the second-best result is underlined.
All methods are implemented without using exemplars.
}
\label{tab:main_results}
\small
\setlength{\tabcolsep}{4pt}
\renewcommand{\arraystretch}{1.12}
\resizebox{\textwidth}{!}{%
\begin{tabular}{lcccccccc}
\toprule
\multirow{2}{*}{Methods}
& \multirow{2}{*}{+Params (M)}
& \multirow{2}{*}{+FLOPs (G)}
& \multicolumn{2}{c}{ImageNet-R}
& \multicolumn{2}{c}{Omnibenchmark}
& \multicolumn{2}{c}{DomainNet} \\
\cmidrule(lr){4-5} \cmidrule(lr){6-7} \cmidrule(lr){8-9}
& &
& $\bar{A}$ & $A_T$
& $\bar{A}$ & $A_T$
& $\bar{A}$ & $A_T$ \\
\midrule

L2P*
& 0.00 & 0.00
& 66.66$\pm$0.36 & 59.78$\pm$0.56
& 73.01$\pm$0.33 & 63.29$\pm$0.26
& 0.00$\pm$0.00 & 0.00$\pm$0.00 \\

DualPrompt*
& 0.00 & 0.00
& 70.29$\pm$0.25 & 66.50$\pm$0.65
& 77.68$\pm$0.18 & 67.45$\pm$0.45
& 0.00$\pm$0.00 & 0.00$\pm$0.00 \\

CODA-Prompt*
& 0.00 & 0.00
& 76.76$\pm$0.59 & 72.98$\pm$0.47
& 79.18$\pm$0.42 & 70.01$\pm$0.65
& 0.00$\pm$0.00 & 0.00$\pm$0.00 \\

ACIL*
& 0.00 & 0.00
& 75.17$\pm$0.05 & 68.40$\pm$0.22
& 84.26$\pm$0.35 & 76.11$\pm$0.52
& 0.00$\pm$0.00 & 0.00$\pm$0.00 \\

SLCA*
& 0.00 & 0.00
& 82.26$\pm$0.26 & 76.82$\pm$0.34
& 82.46$\pm$0.26 & 74.24$\pm$0.19
& 0.00$\pm$0.00 & 0.00$\pm$0.00 \\

SSIAT
& 0.00 & 0.00
& 83.20 & 78.85
& -- & --
& 0.00$\pm$0.00 & 0.00$\pm$0.00 \\

FeCAM*
& 0.00 & 0.00
& 79.02$\pm$0.22 & 72.53$\pm$0.18
& 83.74$\pm$0.16 & 77.18$\pm$0.22
& 0.00$\pm$0.00 & 0.00$\pm$0.00 \\

RanPAC*
& 0.00 & 0.00
& 82.12$\pm$0.31 & 77.55$\pm$0.16
& 85.04$\pm$0.12 & 78.83$\pm$0.18
& 0.00$\pm$0.00 & 0.00$\pm$0.00 \\

APER*
& 0.00 & 0.00
& 75.45$\pm$0.11 & 67.32$\pm$0.04
& 80.72$\pm$0.20 & 74.34$\pm$0.14
& 0.00$\pm$0.00 & 0.00$\pm$0.00 \\

EASE*
& 0.00 & 0.00
& 81.75$\pm$0.32 & 76.20$\pm$0.19
& 75.71$\pm$0.45 & 68.59$\pm$0.40
& 0.00$\pm$0.00 & 0.00$\pm$0.00 \\

COFiMA*
& 0.00 & 0.00
& 82.05$\pm$0.15 & 76.43$\pm$0.12
& 81.69$\pm$0.40 & 73.38$\pm$0.28
& 0.00$\pm$0.00 & 0.00$\pm$0.00 \\

MOS*
& 0.00 & 0.00
& 82.59$\pm$0.34 & 77.80$\pm$0.20
& \underline{86.10$\pm$0.11} & \underline{80.12$\pm$0.15}
& 0.00$\pm$0.00 & 0.00$\pm$0.00 \\

TUNA
& 0.00 & 0.00
& \underline{84.22} & \underline{79.42}
& -- & --
& 0.00$\pm$0.00 & 0.00$\pm$0.00 \\

\textbf{MiN*}
& \textbf{0.00} & \textbf{0.00}
& \textbf{85.18$\pm$0.19} & \textbf{79.75$\pm$0.07}
& \textbf{87.36$\pm$0.04} & \textbf{80.55$\pm$0.10}
& \textbf{0.00$\pm$0.00} & \textbf{0.00$\pm$0.00} \\

\midrule
\textbf{SPIE (ours)}
& 0.00 & 0.00
& 0.00$\pm$0.00 & 0.00$\pm$0.00
& 0.00$\pm$0.00 & 0.00$\pm$0.00
& 0.00$\pm$0.00 & 0.00$\pm$0.00 \\

\bottomrule
\end{tabular}%
}
\end{table}

% Table 2: CIFAR-100, CUB-200

\begin{table*}[h]
\centering
\caption{Average and last performance comparison on CIFAR-100 and CUB-200 with ViT-B/16-IN21K as the backbone. All methods are implemented without using exemplars.}
\label{tab:cifar_cub}
\small
\setlength{\tabcolsep}{3.5pt}
\renewcommand{\arraystretch}{1.1}
\resizebox{\textwidth}{!}{
\begin{tabular}{lcccccccc}
\toprule
& \multicolumn{4}{c}{CIFAR-100} & \multicolumn{4}{c}{CUB-200} \\
\cmidrule(lr){2-5} \cmidrule(lr){6-9}
Methods & \multicolumn{2}{c}{10 steps} & \multicolumn{2}{c}{20 steps}
        & \multicolumn{2}{c}{10 steps} & \multicolumn{2}{c}{20 steps} \\
\cmidrule(lr){2-3} \cmidrule(lr){4-5}
\cmidrule(lr){6-7} \cmidrule(lr){8-9}
& $\bar{A}$ & $A_T$ & $\bar{A}$ & $A_T$
& $\bar{A}$ & $A_T$ & $\bar{A}$ & $A_T$ \\
\midrule
L2P*
& 85.92$_{\pm 0.78}$ & 79.19$_{\pm 0.57}$
& 81.90$_{\pm 0.54}$ & 73.93$_{\pm 0.33}$
& 84.29$_{\pm 0.32}$ & 74.51$_{\pm 0.36}$
& 81.75$_{\pm 0.18}$ & 70.31$_{\pm 0.25}$ \\
DualPrompt*
& 89.65$_{\pm 0.55}$ & 84.89$_{\pm 0.79}$
& 85.57$_{\pm 0.62}$ & 79.27$_{\pm 0.85}$
& 84.39$_{\pm 0.54}$ & 73.45$_{\pm 0.75}$
& 83.79$_{\pm 0.65}$ & 72.48$_{\pm 1.05}$ \\
CODA-Prompt*
& 91.05$_{\pm 0.32}$ & 86.44$_{\pm 0.87}$
& 87.51$_{\pm 0.70}$ & 80.87$_{\pm 0.25}$
& 84.15$_{\pm 0.87}$ & 74.00$_{\pm 1.05}$
& 83.89$_{\pm 0.37}$ & 71.33$_{\pm 0.55}$ \\
ACIL*
& 91.21$_{\pm 0.03}$ & 86.74$_{\pm 0.05}$
& 91.34$_{\pm 0.03}$ & 86.78$_{\pm 0.02}$
& 91.74$_{\pm 0.21}$ & 87.22$_{\pm 0.16}$
& 91.83$_{\pm 0.18}$ & 87.09$_{\pm 0.17}$ \\
SLCA*
& 92.67$_{\pm 0.89}$ & 89.30$_{\pm 0.44}$
& \underline{93.32}$_{\pm 0.76}$ & 88.21$_{\pm 0.55}$
& 86.83$_{\pm 0.44}$ & 79.47$_{\pm 0.57}$
& 83.38$_{\pm 0.34}$ & 74.77$_{\pm 0.28}$ \\
FeCAM*
& 93.23$_{\pm 0.15}$ & 89.05$_{\pm 0.17}$
& 91.86$_{\pm 0.11}$ & 87.04$_{\pm 0.26}$
& 92.73$_{\pm 0.27}$ & 88.38$_{\pm 0.33}$
& 92.89$_{\pm 0.25}$ & 88.42$_{\pm 0.20}$ \\
RanPAC*
& 93.25$_{\pm 0.07}$ & 89.55$_{\pm 0.11}$
& 91.81$_{\pm 0.15}$ & 88.69$_{\pm 0.12}$
& \underline{93.30}$_{\pm 0.14}$ & \underline{89.78}$_{\pm 0.15}$
& \underline{93.30}$_{\pm 0.17}$ & \underline{89.27}$_{\pm 0.20}$ \\
APER*
& 92.22$_{\pm 0.05}$ & 87.45$_{\pm 0.14}$
& 90.57$_{\pm 0.09}$ & 85.03$_{\pm 0.11}$
& 88.32$_{\pm 0.15}$ & 85.75$_{\pm 0.10}$
& 88.34$_{\pm 0.08}$ & 85.42$_{\pm 0.05}$ \\
EASE*
& 92.11$_{\pm 0.27}$ & 87.72$_{\pm 0.88}$
& 91.51$_{\pm 0.33}$ & 85.80$_{\pm 0.52}$
& 90.12$_{\pm 0.87}$ & 83.76$_{\pm 1.02}$
& 90.62$_{\pm 0.65}$ & 83.72$_{\pm 0.62}$ \\
COFiMA*
& \underline{93.87}$_{\pm 0.16}$ & 89.77$_{\pm 0.08}$
& 92.86$_{\pm 0.27}$ & 88.09$_{\pm 0.11}$
& 88.17$_{\pm 0.35}$ & 79.64$_{\pm 0.88}$
& 83.52$_{\pm 1.22}$ & 74.77$_{\pm 1.05}$ \\
MOS*
& 93.83$_{\pm 0.12}$ & \underline{90.19}$_{\pm 0.21}$
& 93.10$_{\pm 0.06}$ & \underline{89.10}$_{\pm 0.25}$
& 92.08$_{\pm 0.30}$ & 88.17$_{\pm 0.46}$
& 92.62$_{\pm 0.35}$ & 88.51$_{\pm 0.18}$ \\
MiN*
& \textbf{95.12}$_{\pm 0.05}$ & \textbf{92.12}$_{\pm 0.16}$
& \textbf{94.31}$_{\pm 0.04}$ & \textbf{91.03}$_{\pm 0.29}$
& \textbf{94.00}$_{\pm 0.15}$ & \textbf{91.22}$_{\pm 0.18}$
& \textbf{93.84}$_{\pm 0.12}$ & \textbf{90.54}$_{\pm 0.14}$ \\
\midrule

\textbf{SPIE (ours)}
& -- & --
& -- & --
& -- & --
& -- & -- \\
\bottomrule
\end{tabular}
}
\end{table*}

Tables~\ref{tab:main_results} and~\ref{tab:cifar_cub} report the main results. Several observations are worth noting. 

\textit{[Discussion paragraphs to be filled after results are available. Suggested structure:]}

\textit{Comparison with prompt-based methods.} L2P, DualPrompt, and CODA-Prompt rely on key-query matching to retrieve prompts at test time. As the task horizon lengthens (e.g., 10--20 steps across 200+ classes), prompt confusion accumulates, leading to a notable accuracy gap relative to prototype- and expert-based methods. This gap is most pronounced on DomainNet, where domain shift compounds the retrieval challenge.

\textit{Comparison with prototype-based methods.} RanPAC, FeCAM, and ACIL avoid task inference entirely through nearest-mean classification on frozen features, achieving strong and stable performance across most benchmarks. However, their reliance on a single global feature representation limits their ability to capture task-specific discriminative patterns, as reflected in the plateauing accuracy on fine-grained datasets like CUB-200.

\textit{Comparison with expert-based methods.} MOS, TUNA, and MiN leverage task-adapted representations, yielding the strongest overall performance. However, this comes at a substantially higher parameter budget, as MOS and TUNA train full per-task adapters, while MiN's adaptation cost is also non-trivial. SPIE matches or exceeds these methods while using only scaling vectors per expert, as quantified in Sec.~\ref{sec:efficiency}.

\textit{Cross-dataset trends.} On datasets with large domain gaps (ImageNet-R, DomainNet), structural methods with task-specific adaptation (SLCA, MOS, MiN) consistently outperform one-size-fits-all prototype classifiers. On more homogeneous datasets (CIFAR-100), the performance gap narrows, and even simple prototype methods approach the state-of-the-art. SPIE's soft posterior fusion is designed to benefit most in the former regime, where task ambiguity is high and expert specialization provides the greatest leverage.

\subsection{Ablation Studies}
\label{sec:ablation}

We conduct all ablation experiments on \textbf{ImageNet-R} and \textbf{DomainNet} under the 10-step protocol ($|\mathcal{C}_1|=20$, $|\mathcal{C}_t|=20$, $T=10$). The ablations cover component contributions, expert architecture choices, routing mechanism design, distributional prototype configurations, and posterior fusion behavior. Every design decision in SPIE is isolated and justified.

\subsubsection{Component contribution.}

Table~\ref{tab:component_ablation} decomposes SPIE into a progressive stack. Starting from the simplest baseline (a linear probe on frozen backbone features, SimpleCIL~\cite{zhou2025revisiting}), we incrementally add each component and measure the marginal gain.

\begin{table}[t]
\centering
\caption{Progressive component ablation on ImageNet-R and DomainNet (10-step). Each row adds one component to the baseline above it.}
\label{tab:component_ablation}
\small
\begin{tabular}{lcccc}
\toprule
& \multicolumn{2}{c}{ImageNet-R} & \multicolumn{2}{c}{DomainNet} \\
\cmidrule(lr){2-3} \cmidrule(lr){4-5}
Configuration & $\bar{A}$ & $A_T$ & $\bar{A}$ & $A_T$ \\
\midrule
(a) Frozen backbone + linear probe (SimpleCIL)      & -- & -- & -- & -- \\
 (b) + CosFace class prototypes & -- & -- & -- & -- \\
(c) + Classifier Alignment (CA)                        & -- & -- & -- & -- \\
(d) + Experts, oracle task routing                    & -- & -- & -- & -- \\
(e) + Experts, hard prototype routing                  & -- & -- & -- & -- \\
(f) + Soft posterior fusion \hfill (\textbf{SPIE})    & -- & -- & -- & -- \\
\bottomrule
\end{tabular}
\end{table}

The gap (b)$-$(a) quantifies the benefit of CosFace training over simple linear probing. The gap (c)$-$(b) isolates the contribution of CA in preventing bias toward recent tasks. The gap (e)$-$(c) measures the added value of task-specific expert representations beyond a globally calibrated prototype bank. The gap (d)$-$(e) reveals the penalty of imperfect (hard) task routing versus an oracle. This is the headroom that soft posterior fusion aims to recover. The final gap (f)$-$(e) is the direct contribution of soft posterior inference over hard prototype-based expert selection.

\noindent\textit{Why this table is exhaustive.} It addresses whether gains come from CosFace (b vs a), from CA alone (c vs b), from having multiple experts at all (e vs c), or specifically from soft fusion (f vs e). If (e) already saturates near (d), soft fusion is validated, and if not, it shows headroom remains.

\subsubsection{Expert architecture design.}

\textbf{Adapter type.} We compare three adaptation strategies for the per-task expert, holding all other components (CA, soft posterior, $K=4$) fixed: (i) \emph{no adapter}, with expert tokens plus a local head only, backbone MLPs are frozen, (ii) \emph{full LoRA}~\cite{hu2021lora} adapters with rank $r=16$ per MLP layer (as in TUNA~\cite{wang2025integrating}), which costs $\sim$1.2M learnable parameters per expert, and (iii) our \emph{expert adapter} with learnable scaling vectors ($\sim$52K learnable parameters per expert, $\sim$23\texttimes$ fewer than full LoRA).

\begin{table}[h]
\centering
\caption{Expert adapter type comparison ($K=4$, same routing).}
\label{tab:adapter_type}
\small
\begin{tabular}{lccc}
\toprule
Adapter & +Params/expert (M) & ImageNet-R $\bar{A}$ & DomainNet $\bar{A}$ \\
\midrule
None (tokens + head only) & [X] & -- & -- \\
Full LoRA ($r=16$)        & [X] & -- & -- \\
Frozen-basis (ours)       & [X] & -- & -- \\
\bottomrule
\end{tabular}
\end{table}

This comparison preempts two criticisms: (i) that the gains come merely from adding any low-rank adapter, and (ii) that a full LoRA would perform universally better. The expert adapter is expected to match full LoRA at a fraction of the cost, because the SVD-derived basis already captures the principal directions of the pretrained MLP weights.

\textbf{Expert token count.}
Table~\ref{tab:token_count} evaluates $K \in \{0, 1, 2, 4, 8\}$. $K=0$ corresponds to removing expert tokens entirely and routing only through the backbone CLS feature.

\begin{table}[h]
\centering
\caption{Expert token count $K$ (expert adapter, full SPIE).}
\label{tab:token_count}
\small
\begin{tabular}{lccccc}
\toprule
& \multicolumn{2}{c}{ImageNet-R} & \multicolumn{2}{c}{DomainNet} \\
\cmidrule(lr){2-3} \cmidrule(lr){4-5}
$K$ & $\bar{A}$ & $A_T$ & $\bar{A}$ & $A_T$ \\
\midrule
0 (CLS token only) & -- & -- & -- & -- \\
1                  & -- & -- & -- & -- \\
2                  & -- & -- & -- & -- \\
4 (default)        & -- & -- & -- & -- \\
8                  & -- & -- & -- & -- \\
\bottomrule
\end{tabular}
\end{table}

This addresses whether the expert token mechanism is necessary, or does the adapter alone suffice? The $K=0$ vs $K=4$ gap answers this directly.

\textbf{Basis source.} To verify that the SVD basis from pretrained weights is meaningful (rather than any low-rank perturbation working), we replace the frozen basis matrices $A, B$ (Eq.~\ref{eq:adapt}) with random orthogonal matrices of the same dimensions.

\begin{table}[h]
\centering
\caption{Basis initialization strategy. Both use rank 256, learnable scaling vectors, $K=4$.}
\label{tab:basis_source}
\small
\begin{tabular}{lcc}
\toprule
Basis & ImageNet-R $\bar{A}$ & DomainNet $\bar{A}$ \\
\midrule
Random orthogonal & -- & -- \\
SVD of pretrained weights (ours) & -- & -- \\
\bottomrule
\end{tabular}
\end{table}

This silences the argument that ``any structured low-rank adapter with frozen basis works equally well." If the SVD basis significantly outperforms random, it confirms that inheriting the pretrained weight structure is essential.

\subsubsection{Routing mechanism design.}

\textbf{Routing strategy.}
Table~\ref{tab:routing_strategy} compares five task-routing strategies. Oracle routing provides the upper bound (ground-truth task ID). Uniform weighting ($q(t|x)=1/T$) represents the degenerate case of ignoring task structure. Hard prototype routing uses argmax over the prototype bank. Soft logit-based routing aggregates prototype logits via log-mean-exp~\cite{wang2025integrating}. Soft prototype routing (ours) uses the task-max-proto posterior (Eq.~\ref{eq:task_posterior}).

\begin{table}[t]
\centering
\caption{Routing strategy comparison (full SPIE components, varying only the task posterior $q(t|x)$).}
\label{tab:routing_strategy}
\small
\begin{tabular}{lcccc}
\toprule
& \multicolumn{2}{c}{ImageNet-R} & \multicolumn{2}{c}{DomainNet} \\
\cmidrule(lr){2-3} \cmidrule(lr){4-5}
Routing strategy & $\bar{A}$ & $A_T$ & $\bar{A}$ & $A_T$ \\
\midrule
Uniform $q(t|x) = 1/T$            & -- & -- & -- & -- \\
Hard prototype (argmax)           & -- & -- & -- & -- \\
Soft logit-based (log-mean-exp)   & -- & -- & -- & -- \\
Soft prototype-based (task\_max\_proto, ours) & -- & -- & -- & -- \\
Oracle (ground-truth task)        & -- & -- & -- & -- \\
\bottomrule
\end{tabular}
\end{table}

The uniform-vs-soft gap validates that the posterior carries useful signal (not just expert averaging). The hard-vs-soft gap validates soft fusion itself. The logit-vs-prototype gap isolates the benefit of deriving the routing signal from the calibrated prototype bank rather than raw shared logits.

\textbf{Routing accuracy over time.}
Figure~\ref{fig:routing_accuracy} (to be added) tracks per-task routing accuracy (the fraction of test samples where $\arg\max_t q(t|x)$ matches the true task) as tasks accumulate. We plot three routers: hard prototype routing, soft prototype routing (with argmax for evaluation), and logit-based routing. The purpose is to show that prototype-based routing maintains high fidelity in later tasks, while logit-based routing degrades as the prototype space expands.

\textbf{Does CA matter for routing?}
Table~\ref{tab:ca_routing} isolates whether classifier alignment affects routing quality, or only classification accuracy. We evaluate routing accuracy with and without CA, using the same expert predictions (only the routing signal changes).

\begin{table}[h]
\centering
\caption{Routing accuracy with and without Classifier Alignment (CA). Numbers report the fraction of samples correctly routed to their ground-truth task at $t=10$.}
\label{tab:ca_routing}
\small
\begin{tabular}{lcc}
\toprule
Configuration & ImageNet-R & DomainNet \\
\midrule
SPIE w/o CA (soft routing) & -- & -- \\
SPIE w/ CA (soft routing)  & -- & -- \\
\bottomrule
\end{tabular}
\end{table}

If CA significantly boosts routing accuracy, it confirms that routing and calibration are coupled, as an uncalibrated prototype bank biases the scores toward recent classes, poisoning $q(t|x)$.

\subsubsection{Distributional prototype design.}

\textbf{Covariance mode.}
Table~\ref{tab:cov_mode} compares storage modes for class-conditional feature distributions used in classifier alignment. All experiments use full SPIE ($K=4$, expert adapter, soft posterior routing), varying only the covariance representation.

\begin{table}[t]
\centering
\caption{Covariance storage mode. Storage per class ($d=768$) and accuracy on ImageNet-R and DomainNet (10-step).}
\label{tab:cov_mode}
\small
\begin{tabular}{lcccc}
\toprule
& & \multicolumn{2}{c}{ImageNet-R} & DomainNet \\
\cmidrule(lr){3-4}
Mode & Storage/class & $\bar{A}$ & $A_T$ & $\bar{A}$ \\
\midrule
Mean only (no covariance)   & $d$ (0.8K)      & -- & -- & -- \\
Diagonal variance           & $2d$ (1.5K)     & -- & -- & -- \\
Diag+LowRank ($r=4$)        & $5d$ (3.8K)     & -- & -- & -- \\
Diag+LowRank ($r=8$)        & $9d$ (6.9K)     & -- & -- & -- \\
Diag+LowRank ($r=16$, ours) & $17d$ (13.1K)   & -- & -- & -- \\
Diag+LowRank ($r=32$)       & $33d$ (25.3K)   & -- & -- & -- \\
Full covariance             & $d+d^2$ (590.1K)& -- & -- & -- \\
\bottomrule
\end{tabular}
\end{table}

This addresses (i) whether covariance information is necessary at all (mean-only vs diag+lowrank), (ii) whether simple diagonal variance suffices (diagonal vs lowrank), and (iii) how much rank is needed before saturation. The $r=16$ default is chosen as the knee point. At this rank, accuracy approaches the full-covariance ceiling while storage remains under $3\%$ of full-covariance cost.

\textbf{Classifier alignment epochs and sample budget.}
We ablate the number of CA fine-tuning epochs ($E_{\text{CA}} \in \{10, 20, 30, 50\}$) and the number of synthetic samples generated per class ($N_{\text{sample}} \in \{50, 100, 200, 400\}$) to confirm that the default settings ($E_{\text{CA}}=30$, $N_{\text{sample}} = 160$) are not overtuned.

\subsubsection{Posterior fusion behavior.}

\textbf{Temperature sensitivity.}
Table~\ref{tab:temperature} evaluates the sensitivity of SPIE's final accuracy to the two temperature parameters. We grid-search $\tau_t \in \{0.1, 0.5, 1.0, 2.0, 5.0\}$ (task posterior temperature) and $\tau_e \in \{0.5, 1.0, 2.0, 5.0\}$ (expert softmax temperature), reporting the best, default ($\tau_t\!=\!\tau_e\!=\!1.0$), and worst-case accuracy.

\begin{table}[h]
\centering
\caption{Temperature sensitivity on ImageNet-R (10-step). Default uses $\tau_t=\tau_e=1.0$.}
\label{tab:temperature}
\small
\begin{tabular}{lccc}
\toprule
Configuration & $\bar{A}$ (best) & $\bar{A}$ (default) & $\bar{A}$ (worst) \\
\midrule
Vary $\tau_t$ (fix $\tau_e=1.0$) & -- & -- & -- \\
Vary $\tau_e$ (fix $\tau_t=1.0$) & -- & -- & -- \\
\bottomrule
\end{tabular}
\end{table}

The goal is to show that the default $\tau_t=\tau_e=1.0$ is near-optimal and that SPIE is not brittle to temperature choice. A narrow gap between best and worst demonstrates robustness, while a wide gap would indicate the need for per-dataset tuning (which we do not perform).

\textbf{Per-component accuracy breakdown.}
Table~\ref{tab:per_component} reports the per-task accuracy of each expert on its own task's test data after training (before fusion), the accuracy of the prototype bank on all seen classes, and the accuracy of the full SPIE posterior fusion. This decomposition verifies that no single component dominates. The prototype bank provides a calibrated base. Experts provide task-specialized discrimination. Fusion combines both strengths without either component acting as a bottleneck.

\begin{table}[h]
\centering
\caption{Per-component accuracy breakdown at $t=10$ on ImageNet-R (10-step). Expert accuracy is micro-averaged over all 10 experts tested on their respective task data.}
\label{tab:per_component}
\small
\begin{tabular}{lcc}
\toprule
Component & Local accuracy & Global $\bar{A}$ \\
\midrule
Prototype bank (on all classes)       & -- & -- \\
Experts (avg.\ on own tasks)                                     & -- & -- \\
Expert oracle fusion (upper bound)                               & -- & -- \\
Full SPIE posterior fusion                                       & -- & -- \\
\bottomrule
\end{tabular}
\end{table}

\textbf{Coupled vs.\ decoupled routing.}
A key design claim is that SPIE's routing signal is fully decoupled from expert training. To test this, we implement an alternative where the prototype bank $\mathcal{P}$ used for routing is \emph{not} calibrated by CA. It instead uses the prototype vectors directly after task training, without distributional replay. We compare three conditions, (a) no CA, uncalibrated routing, (b) CA applied, calibrated routing (SPIE), and (c) routing disabled (uniform fusion). The comparison (a) vs (b) vs (c) isolates whether calibration matters for routing specifically, as opposed to merely improving the prototype bank's own predictions.

\subsubsection{Fairness and efficiency.}

\textbf{Parameter-matched comparison.}
A reviewer might argue that SPIE's gains come from its parameter budget rather than its architecture. To preempt this, we increase the adapter size of competing expert-based methods (e.g., TUNA with LoRA rank $r=64$ instead of $r=16$) until their per-task parameter budget matches SPIE's, and report accuracy. Conversely, we reduce SPIE's expert capacity (e.g., $K=1$, lower $r_v$) to match the parameter budget of lighter baselines.

\textbf{Storage-matched comparison.}
Similarly, we allocate SPIE's low-rank covariance storage budget to a nearest competing prototype method (e.g., FeCAM with a matched number of stored feature vectors per class, treated as pseudo-exemplars) and compare performance. This tests whether storing distributional statistics is more effective than storing an equivalent number of raw feature vectors.

\textbf{FLOPs-normalized accuracy.}
We plot accuracy against per-sample inference FLOPs for all structural methods (L2P, DualPrompt, EASE, SLCA, TUNA, MOS, SPIE) on DomainNet at $T=10$, showing that SPIE occupies the Pareto-optimal region, with accuracy competitive with the most expensive methods at a fraction of their FLOPs.

\subsection{Efficiency Analysis}
\label{sec:efficiency}

\noindent\textbf{Parameter growth.}
Table~\ref{tab:param_growth} compares the persistent parameter budget across representative structural CIL methods after 10 incremental tasks on a 200-class dataset. SPIE's per-expert footprint comprises the learnable scaling vectors ($r_v + d_{\text{out}}$ scalars per MLP layer in each block, totaling $\sim$52K per expert) and expert tokens ($K \cdot d = 3.1$K). The frozen SVD basis matrices for VeRA projections are stored once as recomputable projections ($\sim$23.6M) and excluded from persistent storage.

\begin{table}[h]
\centering
\caption{Parameter budget after $T=10$ on a 200-class dataset (ViT-B/16, $d=768$, 12 blocks). Persistent storage excludes recomputable projections.}
\label{tab:param_growth}
\small
\begin{tabular}{lcccc}
\toprule
Method & Learnable/adapter & Expert heads & Statistical storage & Persistent total \\
\midrule
L2P          & [pool]        & [cls]         & 0      & [sum] \\
EASE         & [adapter]     & [cls]         & [cls]  & [sum] \\
SLCA         & [adapter]     & [cls]         & [cls]  & [sum] \\
TUNA         & [loRA]        & [cls]         & [cls]  & [sum] \\
MOS          & [adapter]     & [cls]         & [cls]  & [sum] \\
SPIE (ours)  & [vera]        & [cls]         & [cls]  & \textbf{[sum]} \\
\bottomrule
\end{tabular}
\end{table}

\noindent\textbf{Inference FLOPs.}
We measure per-sample inference FLOPs at the final task using \texttt{fvcore}~\cite{fvcore}. A single frozen ViT-B/16 forward pass costs approximately $17.6$ GFLOPs. As shown in Table~\ref{tab:flops}, methods that sequentially forward through all per-task adapters (MOS, TUNA) scale linearly with the number of tasks, reaching $11\times$ the single-pass cost. SPIE requires only $2\times$. One shared backbone forward handles prototype routing, plus one parallel forward processes all expert tokens (batched multi-expert inference, Sec.~\ref{sec:experts}). This constant $2\times$ overhead is independent of $T$ and is dominated by single-pass methods that avoid per-task forwarding (L2P, DualPrompt, prototype classifiers).

\begin{table}[h]
\centering
\caption{Per-sample inference FLOPs at $T=10$. Baseline: frozen ViT-B/16 single pass ($17.6$ GFLOPs).}
\label{tab:flops}
\small
\begin{tabular}{lcc}
\toprule
Method & Forward passes & Total GFLOPs \\
\midrule
L2P / DualPrompt / CODA-Prompt & $1\times$ & $\sim$17.8 \\
RanPAC / FeCAM / ACIL (prototype) & $1\times$ & 17.6 \\
EASE / SLCA (single adapter selected) & $1\times$ & 17.6 \\
SPIE (ours)                           & $2\times$ & $\sim$35.2 \\
MOS / TUNA                            & $T\times$ & $\sim$193.6 \\
\bottomrule
\end{tabular}
\end{table}

\noindent\textbf{Storage efficiency of low-rank prototypes.}
Table~\ref{tab:storage_proto} compares the per-class statistical storage required by methods that retain feature distribution information. SPIE's diagonal-plus-low-rank representation ($r=16$) stores $d \cdot (r+2)$ scalars per class ($\sim$13.1K at $d=768$). This represents a $45\times$ reduction over full covariance storage ($d+d^2 \approx 590$K) and only a $9\times$ overhead over diagonal-only storage. Across 200 classes, SPIE's total statistical storage is 2.7M scalars.

\begin{table}[h]
\centering
\caption{Statistical storage per class and total for 200 classes ($d=768$, float32).}
\label{tab:storage_proto}
\small
\begin{tabular}{lccc}
\toprule
Method & Storage/class & 200-class total \\
\midrule
FeCAM / SLCA (full cov)       & $d + d^2$ (590.6K)    & 115.3M \\
SPIE (ours, $r=16$)           & $d \cdot (r+2)$ (13.8K) & 2.7M  \\
MOS / TUNA (diagonal)         & $2d$ (1.5K)           & 0.3M  \\
Mean only (no covariance)     & $d$ (0.8K)            & 0.15M \\
\bottomrule
\end{tabular}
\end{table}

\subsection{Additional Analysis}
\label{sec:analysis}

\noindent\textbf{Routing fidelity across tasks.}
Figure~\ref{fig:routing_accuracy} (to be added) plots the task routing accuracy as a function of $t$ for three routers: hard prototype-based, soft prototype-based (ours, evaluated via argmax of $q(t|x)$ for this metric), and logit-based (log-mean-exp). The key finding is that prototype-based routing maintains high fidelity ($>$90\%) throughout all 10 tasks, while logit-based routing degrades monotonically as the prototype space expands and class boundaries become crowded.

\noindent\textbf{Per-step accuracy trajectory.}
Figure~\ref{fig:per_step} (to be added) plots the per-step accuracy $A_t$ for SPIE and the top three competing methods on DomainNet. This reveals \emph{where} SPIE gains accumulate: on early tasks (where the prototype bank alone is sufficient), the gap is small, and on later tasks (where expert specialization is critical and task ambiguity peaks), SPIE pulls ahead.

\noindent\textbf{Effect of task granularity.}
We evaluate SPIE under three step granularities (10-step, 20-step, 50-step) on ImageNet-R and DomainNet. Finer granularity increases the number of experts and the total number of stored prototypes, amplifying the benefit of SPIE's low per-expert parameter cost and compact statistical storage. We report accuracy and total storage at each granularity.

\noindent\textbf{Compatibility with different backbones.}
To verify that SPIE's design is not tied to ViT-B/16 specifically, we report results with ViT-L/14-IN21K on ImageNet-R (10-step). The larger backbone increases the feature dimension $d$, which proportionally increases the per-class storage cost for full-covariance methods (quadratic in $d$) while SPIE's linear-in-$d$ storage keeps the overhead manageable.
